\section{Prior results, failures and assessment}\label{sec:discussion}

\subsection{What was not counted}
Several useful outcomes are outside the twelve-entry tally. Arguments
for F11 and F42 were independently derived and then identified with
prior answers. Achyuth Jayadevan's prior Lean developments for A5 and S5
were reproduced; these remain credited to him. F31's existing
Lei--Zhang counterexample received a
construction audit and remains credited prior work. The
\record{literature/LEDGER.md}{literature ledger} and
\record{reports/RESULT_SCOPE.md}{scope report} retain these distinctions.

N3 produced an obstruction to a step in a published cover construction
and to an attempted repair, together with a restricted argument whose
countable case is prior. It did not settle the general question.
F37 and F39 produced obstructions to algorithmic shortcuts rather than
the requested algorithms. Exact F20 searches and equational attempts
were inconclusive. These failures help explain the research process
without enlarging the success count.

Implementation failures were also kept. A GAP zero exit status can
coexist with an error, so recorded success required inspection of
the output and its mathematical meaning. Negative controls distinguished
rational from integral splitting, matrices from braid equality and
primitive leading pairs from complete integral branches. For B9 an
early return comparison missed a matrix entry and was corrected in
the historical record. The later uniform first-column obstruction is
a different argument. Retaining both makes the correction inspectable.

\subsection{What this experiment measures}
The run produced a portfolio of arguments and implementations under
an explicit wall-clock limit. It demonstrates how an agent can move
between literature, calculations, proof construction and targeted
formalization. It does not measure a probability of solving an arbitrary
open problem. The input is historical, selection was adaptive, prior
knowledge was available and status checks were necessarily incomplete.
No blinded control or complete novelty survey was performed.

There are also different meanings of ``complete''. F34(a) and F38(a)
propose terminating algorithms while importing a full language
construction that was not implemented. N5 and N8 have full candidate
commands, but finite controls cannot establish their completeness
proofs. N9 constructs one fixed ambient group effectively from an
undecidable Diophantine polynomial without printing its entire
numerical presentation. Each is a mathematical claim with a particular
evidence boundary, not a uniform software milestone.

The next assessment should concentrate on the actual mathematical
interfaces, including F38(c)'s shortening argument and N8's full-block
separation, and independently establish attribution. This version
preserves the deadline record for that purpose. Later corrections
should change the current assessment explicitly while leaving the
historical counts and artifacts recoverable.
